% source: An algebraic approach to the Riemann-Roch theorem and the arithmetic theory of function fields (Athanasios Papaioannou), http://www.fen.bilkent.edu.tr/~franz/mat/Riemann-Roch.pdf
% pdf_page: 0003
% retrieved: 2026-07-14
% transcribed_by: page-transcriber (AI vision, no OCR)

\documentclass[11pt]{article}
\usepackage{amsmath,amssymb,amsthm}

% Small-caps theorem style matching the source (DEFINITION/THEOREM headings, italic bodies).
\newtheoremstyle{scplain}{\topsep}{\topsep}{\itshape}{0pt}{\scshape}{.}{ }{\thmname{#1} \thmnote{#3}}
\theoremstyle{scplain}
\newtheorem*{theorem}{Theorem}
\newtheorem*{definition}{Definition}

% Small-caps "Proof." to match the source (standard amsthm definition with \scshape head).
\makeatletter
\renewenvironment{proof}[1][\proofname]{%
  \par\pushQED{\qed}\normalfont\topsep6\p@\@plus6\p@\relax
  \trivlist\item[\hskip\labelsep\scshape #1\@addpunct{.}]\ignorespaces}%
  {\popQED\endtrivlist\@endpefalse}
\makeatother
\renewcommand{\qedsymbol}{$\square$}

% This page uses only standard operators (\min) and standard amssymb symbols
% (\mathbb, \mathfrak, \mathcal, \mathbf, \twoheadrightarrow); no custom operators required.

\begin{document}

\begin{proof}
The claim is obvious in the cases of interest to Number Theory, namely in
Dedekind domains, since Dedekind domain are assumed Noetherian. However, there is
also a general proof. The second part of the proposition,
$[\mathbb{F} : \mathbb{K}(x)] < \infty$, is trivially true since any element
$z \in \mathbb{F}$ is transcendental over $\mathbb{K}$ if and only if
$[\mathbb{F} : \mathbb{K}[z]] < \infty$. For the first part, we need to show that
$x_1, x_2, \ldots, x_n$ are linearly independent over $\mathbb{K}(x)$. Assume the
contrary; then there is a non-trivial relation
\[
  \sum_{i=1}^{n} a_i x_i = 0,
\]
in which the coefficients $a_i$ may be assumed polynomials (otherwise we just
multiply by the least common divisor of their denominators). A further assumption
we can make is that the $a_i$'s are relatively prime; in particular, $x$ does not
divide all of them. If $b_i = a_i(0)$ are the constant terms of the coefficients,
then there is a maximal $j \in \{1, 2, \ldots, n\}$ such that $b_j \neq 0$, but
$b_i = 0$ for all $i > j$. We obtain
\[
  -a_j x_j = \sum_{i \neq j} x_i,
\]
where $a_i \in \mathcal{O}$ (since $x = x_1 \in \mathfrak{m}$) and
$x_i \in x_j \mathfrak{m}$ for $i < j$ by the above conditions; by the
construction of $j$ we also have $a_i = x g_i$ for $i < j$ where $g_i$ is a
polynomial in $x$. By the previous displayed equation we have
\[
  -a_j = \sum_{i<j} a_i \frac{x_i}{x_j} + \sum_{i>j} \frac{x}{x_i} g_i x_i.
\]
All summands on the right hand side belong to $\mathfrak{m}$, and therefore so
does $a_j$. On the other hand, $a_j = b_j + x g_j$ with
$g_j \in \mathbb{K}[x] \subseteq \mathcal{O}$ and $x \in \mathfrak{m}$, so that
$b_j = a_j - x g_j \in \mathfrak{m} \cap \mathbb{K}$, and this is absurd by
theorem 1.1 (b).
\end{proof}

We now proceed with the proof of the claim.

\begin{itemize}
\item[(a)] Assume that $\mathfrak{m}$ is not principal, and let $x_1 = x$ be any
  element of $\mathfrak{m}$. Since $\mathfrak{m} \neq x_1 \mathcal{O}$, there is
  $x_2 \in \mathfrak{m} - x_1 \mathcal{O}$. Since $x_2 x_1^{-1} \notin \mathcal{O}$,
  we must have $x_1 x_2^{-1} \in \mathcal{O}$ or $x_1 \in x_2 \mathcal{O}$.
  Repeating the above argument, we obtain an infinite sequence that satisfies the
  conditions of the above lemma, a contradiction.

\item[(b)] Let $z \in \mathbb{F}$; we may even assume without loss of generality
  that $z \in \mathcal{O}$, since otherwise $z^{-1} \in \mathcal{O}$. If $z$ is a
  unit, then $z = z t^0$, and this representation is unique. Otherwise, there is a
  unique maximal $m \in \mathbb{N}$, such that $z \in t^m \mathcal{O}$, hence
  $z = u t^m$ for some $u \in \mathcal{O}$; this follows by the lemma we showed
  above. Note that $u$ must belong to $\mathcal{O}^\times$, by the construction of
  $m$.

\item[(c)] For this third claim, note that there is $n \in \mathbb{N}$, such that
  $n = \min\{m \mid x = u t^m, x \in \mathfrak{n}\}$. We claim that
  $\mathfrak{n} = t^n \mathcal{O}$, and the proof of this assertion is trivial.
\end{itemize}

\begin{flushright}$\square$\end{flushright}

\medskip
The second class of objects that will be used in the proof of the Riemann-Roch
theorem is \emph{places}. These are defined in the following

\begin{definition}[1.3]
A place $P$ of the function field $\mathbb{F}/\mathbb{K}$ is the maximal ideal of
some valuation ring $\mathcal{O}$ of $\mathbb{F}/\mathbb{K}$. Any generator $t$ of
$P$ is called a parameter for $P$.

\indent The set of all places of $\mathbb{F}/\mathbb{K}$ will be denoted by
$\mathbf{P}_{\mathbb{F}}$.
\end{definition}

As we shall soon see, places induce and are induced by a particular class of
mappings, valuations:

\begin{definition}[1.4]
A valuation of a function field $\mathbb{F}/\mathbb{K}$ is a surjective
homomorphism $v : \mathbb{F}^* \twoheadrightarrow \mathbb{Z}$, where $\mathbb{F}^*$
denotes the multiplicative group of nonzero elements of the field $\mathbb{F}$,
such that:
\begin{itemize}
\item[(a)] $v\,|_{\mathbb{K}^*} \equiv \mathbf{0}$.

\item[(b)] (\textsc{Triangle Inequality}) $v(x + y) \geq \min\{v(x), v(y)\}$, for
  any pair of $x, y \in \mathbb{F}$.
\end{itemize}
\indent In this context, it is convenient to formally define $v(0) = \infty$; $0$
is the only element of $\mathbb{F}$ mapped to $\infty$.
\end{definition}

\begin{center}3\end{center}

\end{document}
